Federated learning (FL) is a distributed machine learning paradigm that enables multiple clients to collaboratively train models without sharing raw data~\cite{mcmahan2017communication}, making it particularly valuable in privacy-sensitive domains such as healthcare~\cite{rieke2020future}, edge computing~\cite{lim2020federated}, and multi-institutional research~\cite{kairouz2021advances}. As FL deployments grow in scale, they introduce a distinct \emph{observability} challenge that existing tooling does not address. A single training run can involve many geographically distributed clients with heterogeneous data distributions, compute capabilities, communication costs, and failure modes~\cite{li2020federated}. Without a unified monitoring layer, practitioners must manually reconcile disparate logs, re-implement per-client metric schemas for each new project, and hard-code a choice of monitoring backend up front.

General-purpose experiment tracking tools — Weights \& Biases~\cite{wandb}, MLflow~\cite{mlflow}, and TensorBoard~\cite{tensorboard} — are effective for centralized training but have no native notion of a federated round, a client population, per-client staleness, or client geography. Adapting them to FL typically means hand-rolling a per-project naming scheme for per-client metrics and recomputing derived quantities such as gradient divergence and communication volume in every codebase. Modern FL frameworks — APPFL~\cite{ryu2022appfl,li2025advances}, Flower~\cite{beutel2020flower}, NVIDIA FLARE~\cite{roth2022nvidia}, and FedML~\cite{he2020fedml} — each ship their own logging facilities, but these are tightly coupled to the framework's runtime. Switching frameworks means re-learning a different logging surface and rebuilding custom dashboards. Research on heterogeneity and straggler management~\cite{lai2021oort,li2024fedcompass,xie2019asynchronous,nguyen2022federated} defines \emph{what} to measure in FL but does not provide \emph{how to measure it} in a portable way.

HiveWatch addresses this gap. It provides a framework-agnostic runtime API and an FL-native event schema — covering round starts, per-client updates, round summaries, failures, and checkpoints — that sits between any FL framework and any combination of monitoring backends. The same instrumentation code can simultaneously stream experiment data to local replayable artifacts, Weights \& Biases, and MLflow, with derived metrics such as gradient divergence, byte totals, and per-round wall time computed once inside HiveWatch and forwarded to every configured backend. The goal is to separate the \emph{what} of observability (the FL-specific schema) from the \emph{where} (the choice of backend), so that practitioners can instrument once and monitor anywhere.
